\documentclass{standalone}
\usepackage{circuitikz}
\usetikzlibrary{calc}
\definecolor{circuitnotemark}{HTML}{FE00FE}
\begin{document}
\begin{circuitikz}[american, line width=0.8pt]
\ctikzset{bipoles/length=1.2cm}
\coordinate (x1p75y2) at (1.5,-2);
\coordinate (x3p75y2) at (5.5,-2);
\coordinate (x5p75y2) at (9.5,-2);
\coordinate (x7p75y2) at (13.5,-2);
\coordinate (x9p5y2) at (17,-2);
\coordinate (x11p5y2) at (21,-2);
\coordinate (x13p5y2) at (25,-2);
\coordinate (x15p5y2) at (29,-2);
\coordinate (x1p75y4) at (1.5,-6);
\coordinate (x3p75y4) at (5.5,-6);
\coordinate (x5p75y4) at (9.5,-6);
\coordinate (x7p75y4) at (13.5,-6);
\coordinate (x9p75y4) at (17.5,-6);
\coordinate (x11p75y4) at (21.5,-6);
\coordinate (x13p5y4) at (25,-6);
\coordinate (x15p25y4) at (28.5,-6);
\coordinate (x1p75y6) at (1.5,-10);
\coordinate (x3p75y6) at (5.5,-10);
\coordinate (x5p5y6) at (9,-10);
\node[dipchip, num pins=8, hide numbers, font=\scriptsize, /tikz/circuitikz/multipoles/dipchip/width=1.5, /tikz/circuitikz/multipoles/dipchip/pin spacing=0.416667, yshift=-0.25cm] (part-U1) at (x1p75y2) {}; % line 4
\node[font=\scriptsize, anchor=north] at (part-U1.south) {$\mathrm{NE555}$}; % line 4
\node[anchor=south] at (part-U1.north) {$U_{1}$}; % line 4
\node[circuitnotemark, font=\tiny, anchor=west, xshift=2pt, inner sep=0] at (part-U1.bpin 1) {X}; % line 4
\node[circuitnotemark, font=\tiny, anchor=west, xshift=2pt, inner sep=0] at (part-U1.bpin 2) {X}; % line 4
\node[circuitnotemark, font=\tiny, anchor=west, xshift=2pt, inner sep=0] at (part-U1.bpin 3) {X}; % line 4
\node[circuitnotemark, font=\tiny, anchor=west, xshift=2pt, inner sep=0] at (part-U1.bpin 4) {X}; % line 4
\node[circuitnotemark, font=\tiny, anchor=east, xshift=-2pt, inner sep=0] at (part-U1.bpin 5) {X}; % line 4
\node[circuitnotemark, font=\tiny, anchor=east, xshift=-2pt, inner sep=0] at (part-U1.bpin 6) {X}; % line 4
\node[circuitnotemark, font=\tiny, anchor=east, xshift=-2pt, inner sep=0] at (part-U1.bpin 7) {X}; % line 4
\node[circuitnotemark, font=\tiny, anchor=east, xshift=-2pt, inner sep=0] at (part-U1.bpin 8) {X}; % line 4
\node[dipchip, num pins=8, hide numbers, font=\scriptsize, /tikz/circuitikz/multipoles/dipchip/width=1.5, /tikz/circuitikz/multipoles/dipchip/pin spacing=0.416667, yshift=-0.25cm] (part-U2) at (x3p75y2) {}; % line 5
\node[font=\scriptsize, anchor=north] at (part-U2.south) {$\mathrm{TLC555}$}; % line 5
\node[anchor=south] at (part-U2.north) {$U_{2}$}; % line 5
\node[circuitnotemark, font=\tiny, anchor=west, xshift=2pt, inner sep=0] at (part-U2.bpin 1) {X}; % line 5
\node[circuitnotemark, font=\tiny, anchor=west, xshift=2pt, inner sep=0] at (part-U2.bpin 2) {X}; % line 5
\node[circuitnotemark, font=\tiny, anchor=west, xshift=2pt, inner sep=0] at (part-U2.bpin 3) {X}; % line 5
\node[circuitnotemark, font=\tiny, anchor=west, xshift=2pt, inner sep=0] at (part-U2.bpin 4) {X}; % line 5
\node[circuitnotemark, font=\tiny, anchor=east, xshift=-2pt, inner sep=0] at (part-U2.bpin 5) {X}; % line 5
\node[circuitnotemark, font=\tiny, anchor=east, xshift=-2pt, inner sep=0] at (part-U2.bpin 6) {X}; % line 5
\node[circuitnotemark, font=\tiny, anchor=east, xshift=-2pt, inner sep=0] at (part-U2.bpin 7) {X}; % line 5
\node[circuitnotemark, font=\tiny, anchor=east, xshift=-2pt, inner sep=0] at (part-U2.bpin 8) {X}; % line 5
\node[dipchip, num pins=8, hide numbers, font=\scriptsize, /tikz/circuitikz/multipoles/dipchip/width=1.4, /tikz/circuitikz/multipoles/dipchip/pin spacing=0.416667, yshift=-0.25cm] (part-U3) at (x5p75y2) {}; % line 6
\node[font=\scriptsize, anchor=north] at (part-U3.south) {$\mathrm{LM358}$}; % line 6
\node[anchor=south] at (part-U3.north) {$U_{3}$}; % line 6
\node[circuitnotemark, font=\tiny, anchor=west, xshift=2pt, inner sep=0] at (part-U3.bpin 1) {X}; % line 6
\node[circuitnotemark, font=\tiny, anchor=west, xshift=2pt, inner sep=0] at (part-U3.bpin 2) {X}; % line 6
\node[circuitnotemark, font=\tiny, anchor=west, xshift=2pt, inner sep=0] at (part-U3.bpin 3) {X}; % line 6
\node[circuitnotemark, font=\tiny, anchor=west, xshift=2pt, inner sep=0] at (part-U3.bpin 4) {X}; % line 6
\node[circuitnotemark, font=\tiny, anchor=east, xshift=-2pt, inner sep=0] at (part-U3.bpin 5) {X}; % line 6
\node[circuitnotemark, font=\tiny, anchor=east, xshift=-2pt, inner sep=0] at (part-U3.bpin 6) {X}; % line 6
\node[circuitnotemark, font=\tiny, anchor=east, xshift=-2pt, inner sep=0] at (part-U3.bpin 7) {X}; % line 6
\node[circuitnotemark, font=\tiny, anchor=east, xshift=-2pt, inner sep=0] at (part-U3.bpin 8) {X}; % line 6
\node[dipchip, num pins=8, hide numbers, font=\scriptsize, /tikz/circuitikz/multipoles/dipchip/width=1.4, /tikz/circuitikz/multipoles/dipchip/pin spacing=0.416667, yshift=-0.25cm] (part-U4) at (x7p75y2) {}; % line 7
\node[font=\scriptsize, anchor=north] at (part-U4.south) {$\mathrm{TL071}$}; % line 7
\node[anchor=south] at (part-U4.north) {$U_{4}$}; % line 7
\node[circuitnotemark, font=\tiny, anchor=west, xshift=2pt, inner sep=0] at (part-U4.bpin 1) {X}; % line 7
\node[circuitnotemark, font=\tiny, anchor=west, xshift=2pt, inner sep=0] at (part-U4.bpin 2) {X}; % line 7
\node[circuitnotemark, font=\tiny, anchor=west, xshift=2pt, inner sep=0] at (part-U4.bpin 3) {X}; % line 7
\node[circuitnotemark, font=\tiny, anchor=west, xshift=2pt, inner sep=0] at (part-U4.bpin 4) {X}; % line 7
\node[circuitnotemark, font=\tiny, anchor=east, xshift=-2pt, inner sep=0] at (part-U4.bpin 5) {X}; % line 7
\node[circuitnotemark, font=\tiny, anchor=east, xshift=-2pt, inner sep=0] at (part-U4.bpin 6) {X}; % line 7
\node[circuitnotemark, font=\tiny, anchor=east, xshift=-2pt, inner sep=0] at (part-U4.bpin 7) {X}; % line 7
\node[circuitnotemark, font=\tiny, anchor=east, xshift=-2pt, inner sep=0] at (part-U4.bpin 8) {X}; % line 7
\node[dipchip, num pins=8, hide numbers, font=\scriptsize, /tikz/circuitikz/multipoles/dipchip/width=1.4, /tikz/circuitikz/multipoles/dipchip/pin spacing=0.416667, yshift=-0.25cm] (part-U5) at (x9p5y2) {}; % line 8
\node[font=\scriptsize, anchor=north] at (part-U5.south) {$\mathrm{TL072}$}; % line 8
\node[anchor=south] at (part-U5.north) {$U_{5}$}; % line 8
\node[circuitnotemark, font=\tiny, anchor=west, xshift=2pt, inner sep=0] at (part-U5.bpin 1) {X}; % line 8
\node[circuitnotemark, font=\tiny, anchor=west, xshift=2pt, inner sep=0] at (part-U5.bpin 2) {X}; % line 8
\node[circuitnotemark, font=\tiny, anchor=west, xshift=2pt, inner sep=0] at (part-U5.bpin 3) {X}; % line 8
\node[circuitnotemark, font=\tiny, anchor=west, xshift=2pt, inner sep=0] at (part-U5.bpin 4) {X}; % line 8
\node[circuitnotemark, font=\tiny, anchor=east, xshift=-2pt, inner sep=0] at (part-U5.bpin 5) {X}; % line 8
\node[circuitnotemark, font=\tiny, anchor=east, xshift=-2pt, inner sep=0] at (part-U5.bpin 6) {X}; % line 8
\node[circuitnotemark, font=\tiny, anchor=east, xshift=-2pt, inner sep=0] at (part-U5.bpin 7) {X}; % line 8
\node[circuitnotemark, font=\tiny, anchor=east, xshift=-2pt, inner sep=0] at (part-U5.bpin 8) {X}; % line 8
\node[dipchip, num pins=8, hide numbers, font=\scriptsize, /tikz/circuitikz/multipoles/dipchip/width=1.4, /tikz/circuitikz/multipoles/dipchip/pin spacing=0.416667, yshift=-0.25cm] (part-U6) at (x11p5y2) {}; % line 9
\node[font=\scriptsize, anchor=north] at (part-U6.south) {$\mathrm{LM393}$}; % line 9
\node[anchor=south] at (part-U6.north) {$U_{6}$}; % line 9
\node[circuitnotemark, font=\tiny, anchor=west, xshift=2pt, inner sep=0] at (part-U6.bpin 1) {X}; % line 9
\node[circuitnotemark, font=\tiny, anchor=west, xshift=2pt, inner sep=0] at (part-U6.bpin 2) {X}; % line 9
\node[circuitnotemark, font=\tiny, anchor=west, xshift=2pt, inner sep=0] at (part-U6.bpin 3) {X}; % line 9
\node[circuitnotemark, font=\tiny, anchor=west, xshift=2pt, inner sep=0] at (part-U6.bpin 4) {X}; % line 9
\node[circuitnotemark, font=\tiny, anchor=east, xshift=-2pt, inner sep=0] at (part-U6.bpin 5) {X}; % line 9
\node[circuitnotemark, font=\tiny, anchor=east, xshift=-2pt, inner sep=0] at (part-U6.bpin 6) {X}; % line 9
\node[circuitnotemark, font=\tiny, anchor=east, xshift=-2pt, inner sep=0] at (part-U6.bpin 7) {X}; % line 9
\node[circuitnotemark, font=\tiny, anchor=east, xshift=-2pt, inner sep=0] at (part-U6.bpin 8) {X}; % line 9
\node[dipchip, num pins=8, hide numbers, font=\scriptsize, /tikz/circuitikz/multipoles/dipchip/width=1.5, /tikz/circuitikz/multipoles/dipchip/pin spacing=0.416667, yshift=-0.25cm] (part-U7) at (x13p5y2) {}; % line 10
\node[font=\scriptsize, anchor=north] at (part-U7.south) {$\mathrm{SA612}$}; % line 10
\node[anchor=south] at (part-U7.north) {$U_{7}$}; % line 10
\node[circuitnotemark, font=\tiny, anchor=west, xshift=2pt, inner sep=0] at (part-U7.bpin 1) {X}; % line 10
\node[circuitnotemark, font=\tiny, anchor=west, xshift=2pt, inner sep=0] at (part-U7.bpin 2) {X}; % line 10
\node[circuitnotemark, font=\tiny, anchor=west, xshift=2pt, inner sep=0] at (part-U7.bpin 3) {X}; % line 10
\node[circuitnotemark, font=\tiny, anchor=west, xshift=2pt, inner sep=0] at (part-U7.bpin 4) {X}; % line 10
\node[circuitnotemark, font=\tiny, anchor=east, xshift=-2pt, inner sep=0] at (part-U7.bpin 5) {X}; % line 10
\node[circuitnotemark, font=\tiny, anchor=east, xshift=-2pt, inner sep=0] at (part-U7.bpin 6) {X}; % line 10
\node[circuitnotemark, font=\tiny, anchor=east, xshift=-2pt, inner sep=0] at (part-U7.bpin 7) {X}; % line 10
\node[circuitnotemark, font=\tiny, anchor=east, xshift=-2pt, inner sep=0] at (part-U7.bpin 8) {X}; % line 10
\node[dipchip, num pins=8, hide numbers, font=\scriptsize, /tikz/circuitikz/multipoles/dipchip/width=1.5, /tikz/circuitikz/multipoles/dipchip/pin spacing=0.416667, yshift=-0.25cm] (part-U8) at (x15p5y2) {}; % line 11
\node[font=\scriptsize, anchor=north] at (part-U8.south) {$\mathrm{MCP6002}$}; % line 11
\node[anchor=south] at (part-U8.north) {$U_{8}$}; % line 11
\node[circuitnotemark, font=\tiny, anchor=west, xshift=2pt, inner sep=0] at (part-U8.bpin 1) {X}; % line 11
\node[circuitnotemark, font=\tiny, anchor=west, xshift=2pt, inner sep=0] at (part-U8.bpin 2) {X}; % line 11
\node[circuitnotemark, font=\tiny, anchor=west, xshift=2pt, inner sep=0] at (part-U8.bpin 3) {X}; % line 11
\node[circuitnotemark, font=\tiny, anchor=west, xshift=2pt, inner sep=0] at (part-U8.bpin 4) {X}; % line 11
\node[circuitnotemark, font=\tiny, anchor=east, xshift=-2pt, inner sep=0] at (part-U8.bpin 5) {X}; % line 11
\node[circuitnotemark, font=\tiny, anchor=east, xshift=-2pt, inner sep=0] at (part-U8.bpin 6) {X}; % line 11
\node[circuitnotemark, font=\tiny, anchor=east, xshift=-2pt, inner sep=0] at (part-U8.bpin 7) {X}; % line 11
\node[circuitnotemark, font=\tiny, anchor=east, xshift=-2pt, inner sep=0] at (part-U8.bpin 8) {X}; % line 11
\node[dipchip, num pins=8, hide numbers, font=\scriptsize, /tikz/circuitikz/multipoles/dipchip/width=1.7, /tikz/circuitikz/multipoles/dipchip/pin spacing=0.416667, yshift=-0.25cm] (part-U9) at (x1p75y4) {}; % line 12
\node[font=\scriptsize, anchor=north] at (part-U9.south) {$\mathrm{LM386}$}; % line 12
\node[anchor=south] at (part-U9.north) {$U_{9}$}; % line 12
\node[circuitnotemark, font=\tiny, anchor=west, xshift=2pt, inner sep=0] at (part-U9.bpin 1) {X}; % line 12
\node[circuitnotemark, font=\tiny, anchor=west, xshift=2pt, inner sep=0] at (part-U9.bpin 2) {X}; % line 12
\node[circuitnotemark, font=\tiny, anchor=west, xshift=2pt, inner sep=0] at (part-U9.bpin 3) {X}; % line 12
\node[circuitnotemark, font=\tiny, anchor=west, xshift=2pt, inner sep=0] at (part-U9.bpin 4) {X}; % line 12
\node[circuitnotemark, font=\tiny, anchor=east, xshift=-2pt, inner sep=0] at (part-U9.bpin 5) {X}; % line 12
\node[circuitnotemark, font=\tiny, anchor=east, xshift=-2pt, inner sep=0] at (part-U9.bpin 6) {X}; % line 12
\node[circuitnotemark, font=\tiny, anchor=east, xshift=-2pt, inner sep=0] at (part-U9.bpin 7) {X}; % line 12
\node[circuitnotemark, font=\tiny, anchor=east, xshift=-2pt, inner sep=0] at (part-U9.bpin 8) {X}; % line 12
\node[dipchip, num pins=8, hide numbers, font=\scriptsize, /tikz/circuitikz/multipoles/dipchip/width=1.4, /tikz/circuitikz/multipoles/dipchip/pin spacing=0.416667, yshift=-0.25cm] (part-U10) at (x3p75y4) {}; % line 13
\node[font=\scriptsize, anchor=north] at (part-U10.south) {$\mathrm{NE5532}$}; % line 13
\node[anchor=south] at (part-U10.north) {$U_{10}$}; % line 13
\node[circuitnotemark, font=\tiny, anchor=west, xshift=2pt, inner sep=0] at (part-U10.bpin 1) {X}; % line 13
\node[circuitnotemark, font=\tiny, anchor=west, xshift=2pt, inner sep=0] at (part-U10.bpin 2) {X}; % line 13
\node[circuitnotemark, font=\tiny, anchor=west, xshift=2pt, inner sep=0] at (part-U10.bpin 3) {X}; % line 13
\node[circuitnotemark, font=\tiny, anchor=west, xshift=2pt, inner sep=0] at (part-U10.bpin 4) {X}; % line 13
\node[circuitnotemark, font=\tiny, anchor=east, xshift=-2pt, inner sep=0] at (part-U10.bpin 5) {X}; % line 13
\node[circuitnotemark, font=\tiny, anchor=east, xshift=-2pt, inner sep=0] at (part-U10.bpin 6) {X}; % line 13
\node[circuitnotemark, font=\tiny, anchor=east, xshift=-2pt, inner sep=0] at (part-U10.bpin 7) {X}; % line 13
\node[circuitnotemark, font=\tiny, anchor=east, xshift=-2pt, inner sep=0] at (part-U10.bpin 8) {X}; % line 13
\node[dipchip, num pins=8, hide numbers, font=\scriptsize, /tikz/circuitikz/multipoles/dipchip/width=1.4, /tikz/circuitikz/multipoles/dipchip/pin spacing=0.416667, yshift=-0.25cm] (part-U11) at (x5p75y4) {}; % line 14
\node[font=\scriptsize, anchor=north] at (part-U11.south) {$\mathrm{NJM4556A}$}; % line 14
\node[anchor=south] at (part-U11.north) {$U_{11}$}; % line 14
\node[circuitnotemark, font=\tiny, anchor=west, xshift=2pt, inner sep=0] at (part-U11.bpin 1) {X}; % line 14
\node[circuitnotemark, font=\tiny, anchor=west, xshift=2pt, inner sep=0] at (part-U11.bpin 2) {X}; % line 14
\node[circuitnotemark, font=\tiny, anchor=west, xshift=2pt, inner sep=0] at (part-U11.bpin 3) {X}; % line 14
\node[circuitnotemark, font=\tiny, anchor=west, xshift=2pt, inner sep=0] at (part-U11.bpin 4) {X}; % line 14
\node[circuitnotemark, font=\tiny, anchor=east, xshift=-2pt, inner sep=0] at (part-U11.bpin 5) {X}; % line 14
\node[circuitnotemark, font=\tiny, anchor=east, xshift=-2pt, inner sep=0] at (part-U11.bpin 6) {X}; % line 14
\node[circuitnotemark, font=\tiny, anchor=east, xshift=-2pt, inner sep=0] at (part-U11.bpin 7) {X}; % line 14
\node[circuitnotemark, font=\tiny, anchor=east, xshift=-2pt, inner sep=0] at (part-U11.bpin 8) {X}; % line 14
\node[dipchip, num pins=8, hide numbers, font=\scriptsize, /tikz/circuitikz/multipoles/dipchip/width=1.9, /tikz/circuitikz/multipoles/dipchip/pin spacing=0.416667, yshift=-0.25cm] (part-U12) at (x7p75y4) {}; % line 15
\node[font=\scriptsize, anchor=north] at (part-U12.south) {$\mathrm{LM741}$}; % line 15
\node[anchor=south] at (part-U12.north) {$U_{12}$}; % line 15
\node[circuitnotemark, font=\tiny, anchor=west, xshift=2pt, inner sep=0] at (part-U12.bpin 1) {X}; % line 15
\node[circuitnotemark, font=\tiny, anchor=west, xshift=2pt, inner sep=0] at (part-U12.bpin 2) {X}; % line 15
\node[circuitnotemark, font=\tiny, anchor=west, xshift=2pt, inner sep=0] at (part-U12.bpin 3) {X}; % line 15
\node[circuitnotemark, font=\tiny, anchor=west, xshift=2pt, inner sep=0] at (part-U12.bpin 4) {X}; % line 15
\node[circuitnotemark, font=\tiny, anchor=east, xshift=-2pt, inner sep=0] at (part-U12.bpin 5) {X}; % line 15
\node[circuitnotemark, font=\tiny, anchor=east, xshift=-2pt, inner sep=0] at (part-U12.bpin 6) {X}; % line 15
\node[circuitnotemark, font=\tiny, anchor=east, xshift=-2pt, inner sep=0] at (part-U12.bpin 7) {X}; % line 15
\node[circuitnotemark, font=\tiny, anchor=east, xshift=-2pt, inner sep=0] at (part-U12.bpin 8) {X}; % line 15
\node[dipchip, num pins=8, hide numbers, font=\scriptsize, /tikz/circuitikz/multipoles/dipchip/width=1.2, /tikz/circuitikz/multipoles/dipchip/pin spacing=0.416667, yshift=-0.25cm] (part-U13) at (x9p75y4) {}; % line 16
\node[font=\scriptsize, anchor=north] at (part-U13.south) {$\mathrm{ATtiny85}$}; % line 16
\node[anchor=south] at (part-U13.north) {$U_{13}$}; % line 16
\node[circuitnotemark, font=\tiny, anchor=west, xshift=2pt, inner sep=0] at (part-U13.bpin 1) {X}; % line 16
\node[circuitnotemark, font=\tiny, anchor=west, xshift=2pt, inner sep=0] at (part-U13.bpin 2) {X}; % line 16
\node[circuitnotemark, font=\tiny, anchor=west, xshift=2pt, inner sep=0] at (part-U13.bpin 3) {X}; % line 16
\node[circuitnotemark, font=\tiny, anchor=west, xshift=2pt, inner sep=0] at (part-U13.bpin 4) {X}; % line 16
\node[circuitnotemark, font=\tiny, anchor=east, xshift=-2pt, inner sep=0] at (part-U13.bpin 5) {X}; % line 16
\node[circuitnotemark, font=\tiny, anchor=east, xshift=-2pt, inner sep=0] at (part-U13.bpin 6) {X}; % line 16
\node[circuitnotemark, font=\tiny, anchor=east, xshift=-2pt, inner sep=0] at (part-U13.bpin 7) {X}; % line 16
\node[circuitnotemark, font=\tiny, anchor=east, xshift=-2pt, inner sep=0] at (part-U13.bpin 8) {X}; % line 16
\node[dipchip, num pins=8, hide numbers, font=\scriptsize, /tikz/circuitikz/multipoles/dipchip/width=1.2, /tikz/circuitikz/multipoles/dipchip/pin spacing=0.416667, yshift=-0.25cm] (part-U14) at (x11p75y4) {}; % line 17
\node[font=\scriptsize, anchor=north] at (part-U14.south) {$\mathrm{PIC12F683}$}; % line 17
\node[anchor=south] at (part-U14.north) {$U_{14}$}; % line 17
\node[circuitnotemark, font=\tiny, anchor=west, xshift=2pt, inner sep=0] at (part-U14.bpin 1) {X}; % line 17
\node[circuitnotemark, font=\tiny, anchor=west, xshift=2pt, inner sep=0] at (part-U14.bpin 2) {X}; % line 17
\node[circuitnotemark, font=\tiny, anchor=west, xshift=2pt, inner sep=0] at (part-U14.bpin 3) {X}; % line 17
\node[circuitnotemark, font=\tiny, anchor=west, xshift=2pt, inner sep=0] at (part-U14.bpin 4) {X}; % line 17
\node[circuitnotemark, font=\tiny, anchor=east, xshift=-2pt, inner sep=0] at (part-U14.bpin 5) {X}; % line 17
\node[circuitnotemark, font=\tiny, anchor=east, xshift=-2pt, inner sep=0] at (part-U14.bpin 6) {X}; % line 17
\node[circuitnotemark, font=\tiny, anchor=east, xshift=-2pt, inner sep=0] at (part-U14.bpin 7) {X}; % line 17
\node[circuitnotemark, font=\tiny, anchor=east, xshift=-2pt, inner sep=0] at (part-U14.bpin 8) {X}; % line 17
\node[dipchip, num pins=8, hide numbers, font=\scriptsize, /tikz/circuitikz/multipoles/dipchip/width=1.2, /tikz/circuitikz/multipoles/dipchip/pin spacing=0.416667, yshift=-0.25cm] (part-U15) at (x13p5y4) {}; % line 18
\node[font=\scriptsize, anchor=north] at (part-U15.south) {$\mathrm{PIC12F1822}$}; % line 18
\node[anchor=south] at (part-U15.north) {$U_{15}$}; % line 18
\node[circuitnotemark, font=\tiny, anchor=west, xshift=2pt, inner sep=0] at (part-U15.bpin 1) {X}; % line 18
\node[circuitnotemark, font=\tiny, anchor=west, xshift=2pt, inner sep=0] at (part-U15.bpin 2) {X}; % line 18
\node[circuitnotemark, font=\tiny, anchor=west, xshift=2pt, inner sep=0] at (part-U15.bpin 3) {X}; % line 18
\node[circuitnotemark, font=\tiny, anchor=west, xshift=2pt, inner sep=0] at (part-U15.bpin 4) {X}; % line 18
\node[circuitnotemark, font=\tiny, anchor=east, xshift=-2pt, inner sep=0] at (part-U15.bpin 5) {X}; % line 18
\node[circuitnotemark, font=\tiny, anchor=east, xshift=-2pt, inner sep=0] at (part-U15.bpin 6) {X}; % line 18
\node[circuitnotemark, font=\tiny, anchor=east, xshift=-2pt, inner sep=0] at (part-U15.bpin 7) {X}; % line 18
\node[circuitnotemark, font=\tiny, anchor=east, xshift=-2pt, inner sep=0] at (part-U15.bpin 8) {X}; % line 18
\node[dipchip, num pins=4, hide numbers, font=\scriptsize, /tikz/circuitikz/multipoles/dipchip/width=1.2, /tikz/circuitikz/multipoles/dipchip/pin spacing=0.416667, yshift=-0.25cm] (part-U16) at (x15p25y4) {}; % line 19
\node[font=\scriptsize, anchor=north] at (part-U16.south) {$\mathrm{3SK291}$}; % line 19
\node[anchor=south] at (part-U16.north) {$U_{16}$}; % line 19
\node[circuitnotemark, font=\tiny, anchor=west, xshift=2pt, inner sep=0] at (part-U16.bpin 1) {X}; % line 19
\node[circuitnotemark, font=\tiny, anchor=west, xshift=2pt, inner sep=0] at (part-U16.bpin 2) {X}; % line 19
\node[circuitnotemark, font=\tiny, anchor=east, xshift=-2pt, inner sep=0] at (part-U16.bpin 3) {X}; % line 19
\node[circuitnotemark, font=\tiny, anchor=east, xshift=-2pt, inner sep=0] at (part-U16.bpin 4) {X}; % line 19
\node[dipchip, num pins=8, hide numbers, font=\scriptsize, /tikz/circuitikz/multipoles/dipchip/width=1.4, /tikz/circuitikz/multipoles/dipchip/pin spacing=0.416667, yshift=-0.25cm] (part-U17) at (x1p75y6) {}; % line 20
\node[font=\scriptsize, anchor=north] at (part-U17.south) {$\mathrm{AE\mbox{-}SI5351A\mbox{-}B}$}; % line 20
\node[anchor=south] at (part-U17.north) {$U_{17}$}; % line 20
\node[circuitnotemark, font=\tiny, anchor=west, xshift=2pt, inner sep=0] at (part-U17.bpin 1) {X}; % line 20
\node[circuitnotemark, font=\tiny, anchor=west, xshift=2pt, inner sep=0] at (part-U17.bpin 2) {X}; % line 20
\node[circuitnotemark, font=\tiny, anchor=west, xshift=2pt, inner sep=0] at (part-U17.bpin 3) {X}; % line 20
\node[circuitnotemark, font=\tiny, anchor=west, xshift=2pt, inner sep=0] at (part-U17.bpin 4) {X}; % line 20
\node[circuitnotemark, font=\tiny, anchor=east, xshift=-2pt, inner sep=0] at (part-U17.bpin 5) {X}; % line 20
\node[circuitnotemark, font=\tiny, anchor=east, xshift=-2pt, inner sep=0] at (part-U17.bpin 6) {X}; % line 20
\node[circuitnotemark, font=\tiny, anchor=east, xshift=-2pt, inner sep=0] at (part-U17.bpin 7) {X}; % line 20
\node[circuitnotemark, font=\tiny, anchor=east, xshift=-2pt, inner sep=0] at (part-U17.bpin 8) {X}; % line 20
\node[dipchip, num pins=8, hide numbers, font=\scriptsize, /tikz/circuitikz/multipoles/dipchip/width=1.2, /tikz/circuitikz/multipoles/dipchip/pin spacing=0.416667, yshift=-0.25cm] (part-U18) at (x3p75y6) {}; % line 21
\node[font=\scriptsize, anchor=north] at (part-U18.south) {$\mathrm{LM75}$}; % line 21
\node[anchor=south] at (part-U18.north) {$U_{18}$}; % line 21
\node[circuitnotemark, font=\tiny, anchor=west, xshift=2pt, inner sep=0] at (part-U18.bpin 1) {X}; % line 21
\node[circuitnotemark, font=\tiny, anchor=west, xshift=2pt, inner sep=0] at (part-U18.bpin 2) {X}; % line 21
\node[circuitnotemark, font=\tiny, anchor=west, xshift=2pt, inner sep=0] at (part-U18.bpin 3) {X}; % line 21
\node[circuitnotemark, font=\tiny, anchor=west, xshift=2pt, inner sep=0] at (part-U18.bpin 4) {X}; % line 21
\node[circuitnotemark, font=\tiny, anchor=east, xshift=-2pt, inner sep=0] at (part-U18.bpin 5) {X}; % line 21
\node[circuitnotemark, font=\tiny, anchor=east, xshift=-2pt, inner sep=0] at (part-U18.bpin 6) {X}; % line 21
\node[circuitnotemark, font=\tiny, anchor=east, xshift=-2pt, inner sep=0] at (part-U18.bpin 7) {X}; % line 21
\node[circuitnotemark, font=\tiny, anchor=east, xshift=-2pt, inner sep=0] at (part-U18.bpin 8) {X}; % line 21
\node[dipchip, num pins=8, hide numbers, font=\scriptsize, /tikz/circuitikz/multipoles/dipchip/width=1.5, /tikz/circuitikz/multipoles/dipchip/pin spacing=0.416667, yshift=-0.25cm] (part-U19) at (x5p5y6) {}; % line 22
\node[font=\scriptsize, anchor=north] at (part-U19.south) {$\mathrm{MCP9808}$}; % line 22
\node[anchor=south] at (part-U19.north) {$U_{19}$}; % line 22
\node[circuitnotemark, font=\tiny, anchor=west, xshift=2pt, inner sep=0] at (part-U19.bpin 1) {X}; % line 22
\node[circuitnotemark, font=\tiny, anchor=west, xshift=2pt, inner sep=0] at (part-U19.bpin 2) {X}; % line 22
\node[circuitnotemark, font=\tiny, anchor=west, xshift=2pt, inner sep=0] at (part-U19.bpin 3) {X}; % line 22
\node[circuitnotemark, font=\tiny, anchor=west, xshift=2pt, inner sep=0] at (part-U19.bpin 4) {X}; % line 22
\node[circuitnotemark, font=\tiny, anchor=east, xshift=-2pt, inner sep=0] at (part-U19.bpin 5) {X}; % line 22
\node[circuitnotemark, font=\tiny, anchor=east, xshift=-2pt, inner sep=0] at (part-U19.bpin 6) {X}; % line 22
\node[circuitnotemark, font=\tiny, anchor=east, xshift=-2pt, inner sep=0] at (part-U19.bpin 7) {X}; % line 22
\node[circuitnotemark, font=\tiny, anchor=east, xshift=-2pt, inner sep=0] at (part-U19.bpin 8) {X}; % line 22
\node[anchor=south west, inner sep=2pt, yshift=4pt, circuitnotemark, font=\LARGE\bfseries] (circuittitle) at (current bounding box.north west) {X};
\path (circuittitle.south west) rectangle ++(13.929,0.853);
\node[anchor=north east, inner sep=2pt, circuitnotemark, font=\large] (circuitstamp) at ([yshift=-0.593cm]current bounding box.south east) {X};
\path (circuitstamp.north east) rectangle ++(-5.08,-0.593);
\end{circuitikz}
\end{document}
